\documentclass[12pt,a4paper]{article}
\usepackage[UTF8,fontset=none]{ctex}  % 中文支持，不自动加载字体
\setCJKmainfont{SimSun}               % 使用系统字体（宋体）
\usepackage{amsmath,amssymb,amsthm}  % 数学公式
\usepackage{graphicx}                % 插入图片
\usepackage{booktabs}                 % 表格
\usepackage{algorithm}                % 算法
\usepackage{algorithmic}              % 算法伪代码
\usepackage{listings}                 % 代码
\usepackage{xcolor}                   % 颜色
\usepackage{hyperref}                 % 超链接
\usepackage{float}                    % 浮动体控制

% 超链接设置
\hypersetup{
    colorlinks=true,
    linkcolor=blue,
    citecolor=green
}

\title{LaTeX高级功能测试}
\author{测试}
\date{\today}

\begin{document}

\maketitle

\tableofcontents  % 目录
\newpage

\section{数学公式测试}

\subsection{行内公式}

爱因斯坦质能方程：$E = mc^2$，其中$E$是能量，$m$是质量，$c$是光速。

\subsection{独立公式}

高斯积分：
\begin{equation}
\int_{-\infty}^{\infty} e^{-x^2} dx = \sqrt{\pi}
\label{eq:gaussian}
\end{equation}

矩阵：
\begin{equation}
\mathbf{A} = \begin{pmatrix}
a_{11} & a_{12} \\
a_{21} & a_{22}
\end{pmatrix}
\label{eq:matrix}
\end{equation}

\subsection{多行公式}

\begin{align}
f(x) &= x^2 + 2x + 1 \\
     &= (x+1)^2 \\
     &= (x+1)(x+1)
\label{eq:align}
\end{align}

\section{表格测试}

\begin{table}[H]
\centering
\caption{测试表格}
\label{tab:test}
\begin{tabular}{lcc}
\toprule
项目 & 数值1 & 数值2 \\
\midrule
项目A & 100 & 200 \\
项目B & 150 & 250 \\
项目C & 200 & 300 \\
\bottomrule
\end{tabular}
\end{table}

\section{图片测试}

如果项目中有图片文件，可以这样插入：

\begin{figure}[H]
\centering
\includegraphics[width=0.6\textwidth]{figures/phase_transition.png}
\caption{逾渗相变图（如果文件存在）}
\label{fig:phase}
\end{figure}

\textbf{注意：}如果图片文件不存在，编译会报错，但这是正常的。

\section{算法测试}

\begin{algorithm}[H]
\caption{测试算法}
\begin{algorithmic}[1]
\REQUIRE 输入参数 $x$
\ENSURE 输出结果 $y$
\STATE 初始化 $y = 0$
\FOR{$i = 1$ to $n$}
    \STATE $y = y + x_i$
\ENDFOR
\RETURN $y$
\end{algorithmic}
\end{algorithm}

\section{代码测试}

Python代码示例：

\begin{lstlisting}[language=Python, caption=Python代码示例]
def hello_world():
    print("Hello, LaTeX!")
    return True

if __name__ == "__main__":
    hello_world()
\end{lstlisting}

\section{引用测试}

引用公式\eqref{eq:gaussian}和表格\ref{tab:test}。

引用图片\ref{fig:phase}（如果图片存在）。

\section{列表测试}

\subsection{无序列表}

\begin{itemize}
    \item 第一项
    \item 第二项
    \begin{itemize}
        \item 子项1
        \item 子项2
    \end{itemize}
    \item 第三项
\end{itemize}

\subsection{有序列表}

\begin{enumerate}
    \item 第一步
    \item 第二步
    \begin{enumerate}
        \item 子步骤1
        \item 子步骤2
    \end{enumerate}
    \item 第三步
\end{enumerate}

\section{定理环境测试}

\begin{theorem}
这是一个测试定理。
\end{theorem}

\begin{proof}
这是证明过程。
\end{proof}

\section{中文测试}

这是一个包含中文的段落。LaTeX可以很好地处理中文内容，包括各种标点符号：，。！？；：""''（）【】。

\section{超链接测试}

访问 \href{https://www.latex-project.org/}{LaTeX官网}。

\section{测试完成}

如果这个PDF能正常生成，并且所有内容都正确显示，说明：

\begin{itemize}
    \item ✅ LaTeX环境配置成功
    \item ✅ 中文支持正常
    \item ✅ 数学公式正常
    \item ✅ 表格正常
    \item ✅ 算法环境正常
    \item ✅ 代码高亮正常
    \item ✅ 引用功能正常
    \item ✅ 目录生成正常
\end{itemize}

\textbf{恭喜！LaTeX环境完全配置成功！} 🎉

\end{document}

